\section{Project Outline}

\subsection*{Project Title}
Energy-Efficient UAV Path Planning in IoT Networks: A Deep Reinforcement Learning Aided Approach

\subsection*{Project Summary}
This project develops a protocol-aware Deep Reinforcement Learning framework for optimising
the trajectory of an Unmanned Aerial Vehicle (UAV) acting as a mobile LoRaWAN gateway over
a field of static IoT sensors. The central problem addressed is the ``SF Monoculture''
congestion collapse, wherein standard Adaptive Data Rate mechanisms fail when link qualities
change rapidly due to UAV mobility, causing all sensors to converge on identical Spreading
Factors and collapsing packet delivery rates.

\subsection*{Objectives}
\begin{enumerate}
    \item Develop a high-fidelity, Gymnasium-compatible simulation environment incorporating
          Two-Ray Ground Reflection propagation, log-normal shadowing, EMA-based ADR latency,
          asynchronous duty cycles, and the LoRa Capture Effect.
    \item Formulate the UAV data collection task as a Markov Decision Process with a
          multi-objective reward function balancing energy efficiency, data throughput,
          and Age-of-Information urgency.
    \item Train a Deep Q-Network (DQN) agent using Domain Randomisation and Curriculum
          Learning to produce a single generalisable policy across a wide operational envelope.
    \item Evaluate the DQN agent against two greedy baseline heuristics (Nearest-Sensor
          Greedy and SF-Aware Max-Throughput Greedy) using a rigorous multi-seed evaluation
          with Welch's $t$-test significance testing.
    \item Characterise the Sim-to-Real gap and identify directions for hardware-in-the-loop
          validation and multi-UAV extension.
\end{enumerate}

\subsection*{Deliverables}
\begin{itemize}
    \item Custom Gymnasium environment implementing the full LoRaWAN physics stack.
    \item Trained DQN policy generalising across grid sizes $100\times100$ to
          $1{,}000\times1{,}000$ and sensor counts $N \in \{10, 20, 30, 40\}$.
    \item Evaluation results across 5 seeds and 8 (grid, $N$) conditions with statistical
          significance testing.
    \item This dissertation report.
\end{itemize}

\subsection*{Timeline}
\begin{center}
\begin{tabular}{ll}
\toprule
\textbf{Milestone} & \textbf{Target} \\
\midrule
Simulation environment design and implementation & Semester 1, Week 8 \\
Initial DQN training (fixed scenario)            & Semester 1, Week 12 \\
Domain Randomisation and Curriculum Learning     & Semester 2, Week 4 \\
Multi-seed evaluation and baseline comparison    & Semester 2, Week 8 \\
Dissertation write-up                            & Semester 2, Week 12 \\
Submission                                       & Semester 2, Week 14 \\
\bottomrule
\end{tabular}
\end{center}
